% Part: first-order-logic
% Chapter: completeness
% Section: construction-of-model

\documentclass[../../../include/open-logic-section]{subfiles}

\begin{document}

\iftag{FOL}
      {\olfileid{fol}{com}{mod}}
      {\olfileid{pl}{com}{mod}}

\olsection{प्रतिमानाची रचना}

\begin{explain}
\iftag{FOL}{सध्या आपण~$\eq$ चा विचार करत नाही; म्हणजे~$\eq$ नसलेला
  !!{sentence}s चा सुसंगत संच~$\Gamma$ पूर्ततायोग्य आहे एवढेच आपल्याला
  दाखवायचे आहे. प्रथम~$\Gamma$ चा विस्तार करून सुसंगत, !!{complete}
  आणि संतृप्त असा संच~$\Gamma^*$ मिळवतो. या बाबतीत
  प्रतिमान~$\Struct{M(\Gamma^*)}$ ची व्याख्या सोपी आहे: आपण~$\Lang{L'}$
  मधील बंद पदांचा संच प्रांत म्हणून घेतो. प्रत्येक !!{constant} ला तेच
  स्थिर नेमतो आणि, अधिक व्यापकपणे, प्रत्येक बंद पद~$t$ साठी
  $\Value{t}{M(\Gamma^*)} = t$ याची खात्री करतो. !!{predicate}s ना
  असे विस्तार नेमतो की एखादे आण्विक !!{sentence} हे
  $\Struct{M(\Gamma^*)}$ मध्ये तेव्हा आणि केवळ तेव्हाच सत्य असते, जेव्हा
  ते~$\Gamma^*$ मध्ये असते. त्यामुळे~$\Gamma^*$ मधील सर्व आण्विक
  !!{sentence}s हे~$\Struct{M(\Gamma^*)}$ मध्ये सत्य होतील, हे उघड आहे.
  आपण सुरुवात केलेला~$\Gamma^*$ सुसंगत, संपूर्ण आणि संतृप्त असेल, तर
  उरलेली वाक्येही सत्य असतात.}{आता सर्व $!A \in \Gamma$ सत्य करणारे
  !!a{valuation} परिभाषित करण्यास आपण तयार आहोत. त्यासाठी प्रथम
  लिंडेनबाउमचे पूर्वप्रमेय लावतो: आपल्याला संपूर्ण आणि सुसंगत असा
  $\Gamma^* \supseteq \Gamma$ मिळतो. $\Gamma^*$ मधील
  !!{propositional variable}s कडून~$\pAssign v(\Gamma^*)$ निश्चित करू.}
\end{explain}

\iftag{FOL}{%
\begin{defn}[पद-प्रतिमान]
\ollabel{defn:termmodel}
$\Gamma^*$ हा भाषा~$\Lang L$ मधील !!{sentence}s चा !!a{complete},
सुसंगत आणि संतृप्त संच असू द्या. $\Gamma^*$ चे
\emph{पद-प्रतिमान}~$\Struct M(\Gamma^*)$ ही पुढीलप्रमाणे परिभाषित केलेली
!!{structure} आहे:
\begin{enumerate}
\item !!{domain}~$\Domain{M(\Gamma^*)}$ हा~$\Lang L$ मधील सर्व बंद
  पदांचा संच आहे.
\item !!a{constant}~$c$ चा निर्वचनार्थ स्वतः~$c$ आहे:
  $\Assign{c}{M(\Gamma^*)} = c$.
\item !!{function}~$f$ ला असे फलन नेमले जाते की बंद पदे $t_1$, \dots,
  $t_n$ ही वितर्के दिली असता त्याचे मूल्य बंद पद $f(t_1, \dots, t_n)$
  असते:
\[
\Assign{f}{M(\Gamma^*)}(t_1, \dots, t_n) = f(t_1,\dots, t_n)
\]
\item $R$ हे $n$-स्थानी !!{predicate} असेल, तर
  \[
  \tuple{t_1, \dots,
  t_n} \in \Assign{R}{M(\Gamma^*)} \text{ तेव्हा आणि केवळ तेव्हाच } \Atom{R}{t_1, \dots,
    t_n} \in \Gamma^*.
  \]
\end{enumerate}
\end{defn}
}{%
\begin{defn}
\ollabel{defn:termmodel}
$\Gamma^*$ हा !!{formula}s चा संपूर्ण सुसंगत संच आहे असे समजा. मग आपण
पुढील व्याख्या करतो:
\[
\pAssign v(\Gamma^*)(p) = \begin{cases}
  \True & \text{जर $p \in \Gamma^*$}\\
  \False & \text{जर $p \notin \Gamma^*$}
  \end{cases}
\]
\end{defn}
}

\iftag{FOL}{%
आपल्याकडे खरोखरच $\Value{t}{M(\Gamma^*)} = t$ आहे, हे आता पडताळू.

\begin{lem}
\ollabel{lem:val-in-termmodel} \olref{defn:termmodel} मधील
$\Struct M(\Gamma^*)$ हे पद-प्रतिमान असू द्या; मग
$\Value{t}{M(\Gamma^*)} = t$.
\end{lem}
\begin{proof}
 ही सिद्धता~$t$ वरील विगमनाने मिळते. $t$ हे !!a{constant} असण्याचा
 आधार-प्रसंग पद-प्रतिमानाच्या व्याख्येवरून थेट मिळतो. विगमन-पायरीसाठी
 $t_1, \ldots, t_n$ ही अशी बंद पदे आहेत की
 $\Value{t_i}{M(\Gamma^*)} = t_i$, आणि $f$ हे $n$-स्थानी
 !!{function} आहे, असे समजा. मग
\begin{align*}
\Value{f(t_1,\ldots,t_n)}{M(\Gamma^*)} &= \Assign{f}{M(\Gamma^*)}(\Value{t_1}
 {M(\Gamma^*)},
\ldots, \Value{t_n}{M(\Gamma^*)}) \\
&= \Assign{f}{M(\Gamma^*)}(t_1, \dots, t_n) \\
&= f(t_1,\dots, t_n),
\end{align*}
आणि म्हणून विगमनाने हे प्रत्येक बंद पद~$t$ साठी लागू होते.
\end{proof}
}

\iftag{FOL}{%
\begin{explain}
  \iftag{prvEx}{एखादी !!^a{structure}~$\Struct{M}$ अस्तित्ववाची संख्यकीकृत
    !!{sentence}~$\lexists[x][!A(x)]$ सत्य करू शकते, पण ती सत्य करते
    असा कोणताही दृष्टांत~$!A(t)$ नसेल.}{}
  \iftag{prvAll}{एखादी !!^a{structure}~$\Struct{M}$ सर्ववाची संख्यकीकृत
    !!{sentence}~$\lforall[x][!A(x)]$ चे सर्व दृष्टांत~$!A(t)$ सत्य करू
    शकते, पण~$\lforall[x][!A(x)]$ सत्य करणार नाही.}{} याचे कारण सामान्यतः
  प्रांत~$\Domain{M}$ मधील प्रत्येक !!{element} हे एखाद्या बंद पदाचे
  मूल्य नसते ($\Struct{M}$ बंद पदांनी आच्छादित नसेल). म्हणूनच
  पूर्ति-संबंधाची व्याख्या चर-मूल्यांकनांच्या साहाय्याने केली जाते. पण आपल्या
  पद-प्रतिमान~$\Struct{M(\Gamma^*)}$ साठी त्याची आवश्यकता नसती---कारण
  ते बंद पदांनी आच्छादित आहे. पुढील निष्पत्तीचा आशय हाच आहे.
\end{explain}

\begin{prop}
\ollabel{prop:quant-termmodel}
\olref{defn:termmodel} मधील $\Struct M(\Gamma^*)$ हे पद-प्रतिमान असू द्या.
\begin{tagenumerate}{prvEx,prvAll}
\tagitem{prvEx}{$\Sat{M(\Gamma^*)}{\lexists[x][!A(x)]}$ तेव्हा आणि केवळ तेव्हाच,
  जेव्हा किमान एका बंद पद~$t$ साठी $\Sat{M(\Gamma^*)}{!A(t)}$.}{}
\tagitem{prvAll}{$\Sat{M(\Gamma^*)}{\lforall[x][!A(x)]}$ तेव्हा आणि केवळ तेव्हाच,
  जेव्हा सर्व बंद पदे~$t$ यांच्यासाठी $\Sat{M(\Gamma^*)}{!A(t)}$.}{}
\end{tagenumerate}
\end{prop}

\begin{proof}
\begin{tagenumerate}{prvEx,prvAll}
\tagitem{prvEx}{%
  \iftag{probEx}{सराव.}{\olref[syn][ass]{prop:sat-quant} नुसार
    $\Sat{M(\Gamma^*)}{\lexists[x][!A(x)]}$ तेव्हा आणि केवळ तेव्हाच,
    जेव्हा किमान एका चर-मूल्यांकन~$s$ साठी
    $\Sat{M(\Gamma^*)}{!A(x)}[s]$. $\Domain{M(\Gamma^*)}$ मध्ये
    $\Lang{L}$ मधील बंद पदे असल्यामुळे, हे तेव्हा आणि केवळ तेव्हाच लागू
    होते, जेव्हा असे किमान एक बंद पद~$t$ असते की $s(x) = t$ आणि
    $\Sat{M(\Gamma^*)}{!A(x)}[s]$. \olref[fol][syn][ext]{prop:ext-formulas}
    नुसार, $s(x) = t$ असेल तेथे $\Sat{M(\Gamma^*)}{!A(x)}[s]$ तेव्हा आणि
    केवळ तेव्हाच, जेव्हा $\Sat{M(\Gamma^*)}{!A(t)}[s]$. तसेच
    \olref[fol][syn][ass]{prop:sentence-sat-true} नुसार,
    $!A(t)$ हे वाक्य असल्यामुळे $\Sat{M(\Gamma^*)}{!A(t)}[s]$ तेव्हा आणि
    केवळ तेव्हाच, जेव्हा $\Sat{M(\Gamma^*)}{!A(t)}$.}}{}
\tagitem{prvAll}{%
  \iftag{probAll}{सराव.}{\olref[syn][ass]{prop:sat-quant} नुसार
    $\Sat{M(\Gamma^*)}{\lforall[x][!A(x)]}$ तेव्हा आणि केवळ तेव्हाच,
    जेव्हा प्रत्येक चर-मूल्यांकन $s$ साठी
    $\Sat{M(\Gamma^*)}{!A(x)}[s]$. लक्षात घ्या की
    $\Domain{M(\Gamma^*)}$ मध्ये~$\Lang{L}$ मधील बंद पदे असतात. त्यामुळे
    प्रत्येक बंद पद~$t$ साठी $s(x) = t$ असे चर-मूल्यांकन मिळते, आणि
    कोणत्याही चर-मूल्यांकनासाठी $s(x)$ हे एखादे बंद पद~$t$ असते.
    \olref[fol][syn][ext]{prop:ext-formulas} नुसार, $s(x) = t$ असेल तेथे
    $\Sat{M(\Gamma^*)}{!A(x)}[s]$ तेव्हा आणि केवळ तेव्हाच, जेव्हा
    $\Sat{M(\Gamma^*)}{!A(t)}[s]$. तसेच
    \olref[fol][syn][ass]{prop:sentence-sat-true} नुसार,
    $!A(t)$ हे वाक्य असल्यामुळे $\Sat{M(\Gamma^*)}{!A(t)}[s]$ तेव्हा आणि
    केवळ तेव्हाच, जेव्हा $\Sat{M(\Gamma^*)}{!A(t)}$.}}{}
\end{tagenumerate}
\end{proof}
}{}

\tagprob[FOL]{probEx,probAll}
\begin{prob}
\olref[fol][com][mod]{prop:quant-termmodel} ची सिद्धता पूर्ण करा.
\end{prob}
\tagendprob

\begin{lem}[सत्यता पूर्वप्रमेय]
\ollabel{lem:truth} \iftag{FOL}{$!A$ मध्ये~$\eq$ येत नाही असे समजा. मग}{}
$\iftag{FOL}{\Sat{M(\Gamma^*)}}{\pSat{v(\Gamma^*)}}{!A}$ तेव्हा आणि केवळ तेव्हाच, जेव्हा $!A \in \Gamma^*$.
\end{lem}

\begin{proof}
दोन्ही दिशा एकाच वेळी, $!A$ वरील विगमनाने सिद्ध करू.
\begin{enumerate}
\tagitem{prvFalse}{\indcase{!A}{\lfalse}
  {$\iftag{FOL}{\Sat/{M(\Gamma^*)}{\lfalse}}{\pSat/{v(\Gamma^*)}{\lfalse}}$
    हे पूर्तिच्या व्याख्येनुसार मिळते. दुसरीकडे, $\Gamma^*$ सुसंगत
    असल्यामुळे $\lfalse \notin \Gamma^*$.}}{}

\tagitem{prvTrue}{\indcase{!A}{\ltrue}
  {$\iftag{FOL}{\Sat{M(\Gamma^*)}}{\pSat{v(\Gamma^*)}}{\ltrue}$
    हे पूर्तिच्या व्याख्येनुसार मिळते. दुसरीकडे, $\Gamma^*$ सुसंगत व
    !!{complete} असल्यामुळे आणि $\Gamma^* \Proves \ltrue$ असल्यामुळे
    $\ltrue \in \Gamma^*$.}}{}

\tagitem{FOL}{\indcase{!A}{R(t_1, \dots, t_n)}
  {$\Sat{M(\Gamma^*)}{\Atom{R}{t_1, \dots, t_n}}$ तेव्हा आणि केवळ तेव्हाच,
    जेव्हा $\tuple{t_1, \dots, t_n} \in \Assign{R}{M(\Gamma^*)}$
    (पूर्तिच्या व्याख्येनुसार); आणि हे तेव्हा आणि केवळ तेव्हाच, जेव्हा
    $R(t_1, \dots, t_n) \in \Gamma^*$ ($\Struct M(\Gamma^*)$ च्या
    रचनेनुसार).}}{\indcase{!A}{p}{$\pSat{v(\Gamma^*)}{p}$ तेव्हा आणि
    केवळ तेव्हाच, जेव्हा $\pAssign v(\Gamma^*)(p) = \True$ (पूर्तिच्या
    व्याख्येनुसार); आणि हे तेव्हा आणि केवळ तेव्हाच, जेव्हा
    $p \in \Gamma^*$ ($\pAssign v(\Gamma^*)$ च्या रचनेनुसार).}}

\tagitem{prvNot}{%
  \indcase{!A}{\lnot !B}
    {$\iftag{FOL}{\Sat{M(\Gamma^*)}}{\pSat{v(\Gamma^*)}}{\indfrm}$ तेव्हा
    आणि केवळ तेव्हाच, जेव्हा
    $\iftag{FOL}{\Sat/{M(\Gamma^*)}{!B}}{\pSat/{v(\Gamma^*)}{!B}}$
    (पूर्तिच्या व्याख्येनुसार). विगमन-गृहीतकानुसार,
    $\iftag{FOL}{\Sat/{M(\Gamma^*)}{!B}}{\pSat/{v(\Gamma^*)}{!B}}$ तेव्हा
    आणि केवळ तेव्हाच, जेव्हा $!B \notin \Gamma^*$. $\Gamma^*$ सुसंगत
    आणि !!{complete} असल्यामुळे $!B \notin \Gamma^*$ तेव्हा आणि केवळ
    तेव्हाच, जेव्हा $\lnot !B \in \Gamma^*$.}}{}

\tagitem{prvAnd}{%
\iftag{probAnd}{%
  \indcase!{!A}{!B \land !C}{}}{%
  \indcase{!A}{!B \land
    !C}{$\iftag{FOL}{\Sat{M(\Gamma^*)}}{\pSat{v(\Gamma^*)}}{\indfrm}$
    तेव्हा आणि केवळ तेव्हाच, जेव्हा
    $\iftag{FOL}{\Sat{M(\Gamma^*)}}{\pSat{v(\Gamma^*)}}{!B}$ आणि
    $\iftag{FOL}{\Sat{M(\Gamma^*)}}{\pSat{v(\Gamma^*)}}{!C}$ हे दोन्ही
    लागू होतात (पूर्तिच्या व्याख्येनुसार); आणि हे तेव्हा आणि केवळ तेव्हाच,
    जेव्हा $!B \in \Gamma^*$ आणि $!C \in \Gamma^*$ हे दोन्ही लागू
    होतात (विगमन-गृहीतकानुसार). तसेच
    \olref[ccs]{prop:ccs}\olref[ccs]{prop:ccs-and} नुसार, हे तेव्हा आणि
    केवळ तेव्हाच लागू होते, जेव्हा $(!B \land !C) \in \Gamma^*$.}}}{}

\tagitem{prvOr}{%
\iftag{probOr}{%
  \indcase!{!A}{!B \lor !C}{}}{%
  \indcase{!A}{!B \lor !C}
    {$\iftag{FOL}{\Sat{M(\Gamma^*)}}{\pSat{v(\Gamma^*)}}{\indfrm}$
    तेव्हा आणि केवळ तेव्हाच, जेव्हा
    $\iftag{FOL}{\Sat{M(\Gamma^*)}}{\pSat{v(\Gamma^*)}}{!B}$ किंवा
    $\iftag{FOL}{\Sat{M(\Gamma^*)}}{\pSat{v(\Gamma^*)}}{!C}$
    (पूर्तिच्या व्याख्येनुसार); आणि हे तेव्हा आणि केवळ तेव्हाच, जेव्हा
    $!B \in \Gamma^*$ किंवा $!C \in \Gamma^*$ (विगमन-गृहीतकानुसार).
    \olref[ccs]{prop:ccs}\olref[ccs]{prop:ccs-or} नुसार, हे तेव्हा आणि
    केवळ तेव्हाच लागू होते, जेव्हा $(!B \lor !C) \in \Gamma^*$.}}}{}

\tagitem{prvIf}{%
\iftag{probIf}{%
  \indcase!{!A}{!B \lif !C}{}}{%
  \indcase{!A}{!B \lif !C}
    {$\iftag{FOL}{\Sat{M(\Gamma^*)}}{\pSat{v(\Gamma^*)}}{\indfrm}$
    तेव्हा आणि केवळ तेव्हाच, जेव्हा
    $\iftag{FOL}{\Sat/{M(\Gamma^*)}}{\pSat/{v(\Gamma^*)}}{!B}$ किंवा $\iftag{FOL}{\Sat{M (\Gamma^*)}}{\pSat{v(\Gamma^*)}}{!C}$
    (पूर्तिच्या व्याख्येनुसार); आणि हे तेव्हा आणि केवळ तेव्हाच, जेव्हा
    $!B \notin \Gamma^*$ किंवा $!C \in \Gamma^*$ (विगमन-गृहीतकानुसार).
    \olref[ccs]{prop:ccs}\olref[ccs]{prop:ccs-if} नुसार, हे तेव्हा आणि
    केवळ तेव्हाच लागू होते, जेव्हा $(!B \lif !C) \in \Gamma^*$.}}}{}

\iftag{FOL}{%
\tagitem{prvAll}{%
  \iftag{probAll}{%
    \indcase!{!A}{\lforall[x][!B(x)]}{}}{%
    \indcase{!A}{\lforall[x][!B(x)]}{$\Sat{M(\Gamma^*)}{\indfrm}$ तेव्हा
      आणि केवळ तेव्हाच, जेव्हा सर्व बंद पदे~$t$ यांच्यासाठी
      $\Sat{M(\Gamma^*)}{!B(t)}$ (\olref{prop:quant-termmodel}).
      विगमन-गृहीतकानुसार हे तेव्हा आणि केवळ तेव्हाच लागू होते, जेव्हा सर्व
      बंद पदे~$t$ यांच्यासाठी $!B(t) \in \Gamma^*$; आणि
      \olref[hen]{prop:saturated-instances} नुसार हे पुन्हा तेव्हा आणि
      केवळ तेव्हाच लागू होते, जेव्हा $\lforall[x][!B(x)] \in \Gamma^*$.}}}{}
%READERNOTE{OLCOM-002: गोठवलेल्या इंग्रजी स्रोताच्या सर्ववाची विगमन-प्रसंगातील अंतिम सूत्रात B ऐवजी A आले आहे. मराठीत प्रसंगाशी जुळणारे B पुनर्स्थापित केले असून गोठवलेले इंग्रजी बाइट्स बदललेले नाहीत.}

\tagitem{prvEx}{%
  \iftag{probEx}{%
    \indcase!{!A}{\lexists[x][!B(x)]}{}}{%
    \indcase{!A}{\lexists[x][!B(x)]}{$\Sat{M(\Gamma^*)}{\indfrm}$ तेव्हा
      आणि केवळ तेव्हाच, जेव्हा किमान एका बंद पद~$t$ साठी
      $\Sat{M(\Gamma^*)}{!B(t)}$ (\olref{prop:quant-termmodel}).
      विगमन-गृहीतकानुसार हे तेव्हा आणि केवळ तेव्हाच लागू होते, जेव्हा किमान
      एका बंद पद~$t$ साठी $!B(t) \in \Gamma^*$. तसेच
      \olref[hen]{prop:saturated-instances} नुसार हे पुन्हा तेव्हा आणि
      केवळ तेव्हाच लागू होते, जेव्हा $\lexists[x][!B(x)] \in \Gamma^*$.}}}{}
}{}
%READERNOTE{OLCOM-003: गोठवलेल्या इंग्रजी स्रोताच्या सत्यता पूर्वप्रमेयातील संख्यकीकृत प्रसंगांत चार वेळा पदे म्हटले आहे; संदर्भित निष्पत्ती आणि संतृप्त-दृष्टांत गुणधर्म बंद पदांसाठी आहेत. मराठीत बंद पदे असे स्पष्ट केले असून गोठवलेले इंग्रजी बाइट्स बदललेले नाहीत.}
\end{enumerate}
\end{proof}


\tagprob[FOL]{probOr,probAnd,probIf,probEx,probAll}

\begin{prob}
\olref[fol][com][mod]{lem:truth} ची सिद्धता पूर्ण करा.
\end{prob}

\tagendprob

\tagprob[notFOL]{probOr,probAnd,probIf}

\begin{prob}
\olref[pl][com][mod]{lem:truth} ची सिद्धता पूर्ण करा.
\end{prob}

\tagendprob

\end{document}
